\section{A globally separated interpolation comparator}
\label{sec:comparator}
The next construction bounds the energy of every minimizer. It works
for arbitrary training labels and enforces the margin on the entire cube.

\begin{lemma}[Localized interpolation]\label{lem:comparator}
For $1\le n\le e^{d/25}$, with probability at least $1-e^{-c_Ld}$
over independent uniform inputs, the following holds simultaneously for
every labeling $y_i\in\{-1,1\}$ and every $A>0$.
There is a field $h$ with
\[
 |h(x)|>A\quad\text{for every }x,\qquad
 y_i h(X_i)>A\quad\text{for every }i,
 \qquad \|h\|_K^2\le C_L A^2(d+n d^\beta).
\]
For $n=0$, the same deterministic energy bound holds using a constant field.
\end{lemma}
\begin{proof}
For $z\in\mathcal X$ define the nonnegative bump
\[
 b_z(x)=20\int_{2/5}^{9/20}(1-2p)^{\operatorname{dist}(x,z)}\,dp.
\]
It satisfies $b_z(z)=1$ and $b_z(x)\le(1/5)^{\operatorname{dist}(x,z)}$.
Its degree-$j$ Walsh coefficients, apart from their translation signs,
are
\[
 \widehat b_j=20\int_{2/5}^{9/20}(1-p)^{d-j}p^j\,dp,
 \qquad
 \binom dj\widehat b_j\le\frac{20}{d+1},\qquad
 \widehat b_j\le(3/5)^d.
\]
The middle inequality follows by extending to the complete beta integral.
By Corollary~\ref{cor:lower-spectrum},
$\lambda_j\ge c_Ld^{-1-\beta}/\binom dj$ for $j\le.49d$.
Hence these degrees contribute at most
\[
 C_Ld^{1+\beta}\sum_{j\le.49d}
       \left[\binom dj\widehat b_j\right]^2\le C_Ld^\beta
\]
to $\|b_z\|_K^2$.
For $j>.49d$, the integrand is increasing in $p$ on the integration
interval, so $\widehat b_j\le.55^{d-j}.45^j$.
Let $v=\lambda_{\min}(K)$ be the smallest matrix eigenvalue.
The same corollary gives $v\ge c_L(d\log d)^{-\beta}$ and
$\lambda_j\ge v/2^d$. Thus the remaining energy is at most
\[
 \frac{2^d}{v}\sum_{j>.49d}\binom dj .55^{2(d-j)}.45^{2j}
 \le\operatorname{poly}(d)e^{-c_*d},
\]
where
\[
 c_*=-[\log2+\mathsf h(.49)+1.02\log(.55)+.98\log(.45)]>.0062,
 \qquad \mathsf h(t)=-t\log t-(1-t)\log(1-t).
\]
The exponent decreases beyond $.49d$, so the binomial entropy bound
applies. We have proved $A_d:=\|b_z\|_K^2\le C_Ld^\beta$.

Put $J(z,w)=\langle b_z,b_w\rangle_K$.
Its normalized Fourier eigenvalues are
$\nu_j=\widehat b_j^2/\lambda_j\le .72^d/v$.
Parseval and Cauchy--Schwarz give
\[
 \mathbb E_{X,Y}|J(X,Y)|
 \le\left(\sum_j\binom dj\nu_j^2\right)^{1/2}
 \le\sqrt{\nu_{\max}A_d}.
\]
For every subset $T$ of the sampled indices, including one chosen after
the sample is seen,
\[
 \left\|\sum_{i\in T}b_{X_i}\right\|_K^2
 \le nA_d+\sum_{i\ne j}|J(X_i,X_j)|.
\]
The cross sum has expectation at most
$n(n-1)\sqrt{\nu_{\max}A_d}$.
Markov's inequality at threshold $n d^\beta$ has exponentially small
failure probability for $n\le e^{.04d}$, because
$\tfrac12\log(1/.72)>.16$.
Thus all subsets simultaneously have energy at most $C_Ln d^\beta$.

Also, except with exponentially small probability, every pair of
different sample indices has inputs at Hamming distance greater than
$.2d$. A union bound costs at most
\[
 n^2\exp[-d(\log2-\mathsf h(.2))],
 \qquad \log2-\mathsf h(.2)>.19.
\]
This event excludes repeated inputs.
At any input at most one center is within distance $.1d$; the sum of
all other bumps is at most $n(.2)^{.1d}=o(1)$.
Consequently, for every selected subset of centers, its bump sum is
$1+o(1)$ at a selected center and at most $.2+o(1)$ everywhere else.

For $T=\{i:y_i=-1\}$, take
\[
 h=2A\left(1-2\sum_{i\in T}b_{X_i}\right).
\]
The preceding estimates prove its strict global and training margins
for all sufficiently large $d$.
Finally $\|1\|_K^2=1/\lambda_0\le C_Ld$ by
Corollary~\ref{cor:lower-spectrum}; the squared triangle inequality
proves the asserted energy bound. For $n=0$, use $h=2A$.
\end{proof}
